\section{Coefficient certificates and additional complete degrees}\label{app:tables}
Every table in this appendix was regenerated from the exact data files by
\texttt{code/make\_tables.py}, which recomputed each printed degree-one to
degree-six coefficient by the finite cut formula \eqref{eq:wordcut}, expanded
the displayed brackets \eqref{eq:z1}--\eqref{eq:z6} and
\eqref{eq:c2}--\eqref{eq:c4} by \eqref{eq:bracketexpand}, reconstructed every
Lyndon table below from its brackets, and expanded the tree polynomials of
\cref{thm:tree} into words, comparing all results exactly before writing
the file.

\subsection{All nonzero associative BCH coefficients through degree six}
For each listed word $w$ the table gives $c(w)$ of \eqref{eq:wordcut}; every
unlisted word of degree at most six has coefficient zero, for both the cut
formula and the expansion of \eqref{eq:z1}--\eqref{eq:z6}. Among the
$2+4+8+16+32+64=126$ words there are $72$ nonzero entries. This is a complete
finite certificate for the displayed polynomial through degree six.
\begingroup\small
\begin{longtable}{@{}lrlrlr@{}}
\caption{All nonzero word coefficients $c(w)$ of $Z_1,\ldots,Z_6$.}\label{tab:words6}\\
\toprule
Word $w$ & $c(w)$ & Word $w$ & $c(w)$ & Word $w$ & $c(w)$\\
\midrule\endfirsthead
\multicolumn{6}{c}{\tablename\ \thetable\ (continued)}\\
\toprule
Word $w$ & $c(w)$ & Word $w$ & $c(w)$ & Word $w$ & $c(w)$\\
\midrule\endhead
\bottomrule\endfoot
\multicolumn{6}{@{}l}{\textit{Degree 1}}\\[2pt]
\texttt{X} & $1$ & \texttt{Y} & $1$ &  &  \\
\addlinespace[3pt]
\multicolumn{6}{@{}l}{\textit{Degree 2}}\\[2pt]
\texttt{XY} & $\tfrac{1}{2}$ & \texttt{YX} & $-\tfrac{1}{2}$ &  &  \\
\addlinespace[3pt]
\multicolumn{6}{@{}l}{\textit{Degree 3}}\\[2pt]
\texttt{XXY} & $\tfrac{1}{12}$ & \texttt{XYY} & $\tfrac{1}{12}$ & \texttt{YXY} & $-\tfrac{1}{6}$ \\
\texttt{XYX} & $-\tfrac{1}{6}$ & \texttt{YXX} & $\tfrac{1}{12}$ & \texttt{YYX} & $\tfrac{1}{12}$ \\
\addlinespace[3pt]
\multicolumn{6}{@{}l}{\textit{Degree 4}}\\[2pt]
\texttt{XXYY} & $\tfrac{1}{24}$ & \texttt{YXYX} & $\tfrac{1}{12}$ &  &  \\
\texttt{XYXY} & $-\tfrac{1}{12}$ & \texttt{YYXX} & $-\tfrac{1}{24}$ &  &  \\
\addlinespace[3pt]
\multicolumn{6}{@{}l}{\textit{Degree 5}}\\[2pt]
\texttt{XXXXY} & $-\tfrac{1}{720}$ & \texttt{XYXYY} & $-\tfrac{1}{120}$ & \texttt{YXYXY} & $\tfrac{1}{30}$ \\
\texttt{XXXYX} & $\tfrac{1}{180}$ & \texttt{XYYXX} & $-\tfrac{1}{120}$ & \texttt{YXYYX} & $-\tfrac{1}{120}$ \\
\texttt{XXXYY} & $\tfrac{1}{180}$ & \texttt{XYYXY} & $-\tfrac{1}{120}$ & \texttt{YXYYY} & $\tfrac{1}{180}$ \\
\texttt{XXYXX} & $-\tfrac{1}{120}$ & \texttt{XYYYX} & $\tfrac{1}{180}$ & \texttt{YYXXX} & $\tfrac{1}{180}$ \\
\texttt{XXYXY} & $-\tfrac{1}{120}$ & \texttt{XYYYY} & $-\tfrac{1}{720}$ & \texttt{YYXXY} & $-\tfrac{1}{120}$ \\
\texttt{XXYYX} & $-\tfrac{1}{120}$ & \texttt{YXXXX} & $-\tfrac{1}{720}$ & \texttt{YYXYX} & $-\tfrac{1}{120}$ \\
\texttt{XXYYY} & $\tfrac{1}{180}$ & \texttt{YXXXY} & $\tfrac{1}{180}$ & \texttt{YYXYY} & $-\tfrac{1}{120}$ \\
\texttt{XYXXX} & $\tfrac{1}{180}$ & \texttt{YXXYX} & $-\tfrac{1}{120}$ & \texttt{YYYXX} & $\tfrac{1}{180}$ \\
\texttt{XYXXY} & $-\tfrac{1}{120}$ & \texttt{YXXYY} & $-\tfrac{1}{120}$ & \texttt{YYYXY} & $\tfrac{1}{180}$ \\
\texttt{XYXYX} & $\tfrac{1}{30}$ & \texttt{YXYXX} & $-\tfrac{1}{120}$ & \texttt{YYYYX} & $-\tfrac{1}{720}$ \\
\addlinespace[3pt]
\multicolumn{6}{@{}l}{\textit{Degree 6}}\\[2pt]
\texttt{XXXXYY} & $-\tfrac{1}{1440}$ & \texttt{XYXYYY} & $\tfrac{1}{360}$ & \texttt{YXYYYX} & $-\tfrac{1}{360}$ \\
\texttt{XXXYXY} & $\tfrac{1}{360}$ & \texttt{XYYXXY} & $-\tfrac{1}{240}$ & \texttt{YYXXXX} & $\tfrac{1}{1440}$ \\
\texttt{XXXYYY} & $\tfrac{1}{360}$ & \texttt{XYYXYY} & $-\tfrac{1}{240}$ & \texttt{YYXXYX} & $\tfrac{1}{240}$ \\
\texttt{XXYXXY} & $-\tfrac{1}{240}$ & \texttt{XYYYXY} & $\tfrac{1}{360}$ & \texttt{YYXYXX} & $\tfrac{1}{240}$ \\
\texttt{XXYXYY} & $-\tfrac{1}{240}$ & \texttt{YXXXYX} & $-\tfrac{1}{360}$ & \texttt{YYXYYX} & $\tfrac{1}{240}$ \\
\texttt{XXYYXY} & $-\tfrac{1}{240}$ & \texttt{YXXYXX} & $\tfrac{1}{240}$ & \texttt{YYYXXX} & $-\tfrac{1}{360}$ \\
\texttt{XXYYYY} & $-\tfrac{1}{1440}$ & \texttt{YXXYYX} & $\tfrac{1}{240}$ & \texttt{YYYXYX} & $-\tfrac{1}{360}$ \\
\texttt{XYXXXY} & $\tfrac{1}{360}$ & \texttt{YXYXXX} & $-\tfrac{1}{360}$ & \texttt{YYYYXX} & $\tfrac{1}{1440}$ \\
\texttt{XYXXYY} & $-\tfrac{1}{240}$ & \texttt{YXYXYX} & $-\tfrac{1}{60}$ &  &  \\
\texttt{XYXYXY} & $\tfrac{1}{60}$ & \texttt{YXYYXX} & $\tfrac{1}{240}$ &  &  \\
\end{longtable}\endgroup

\subsection{Complete seventh and eighth BCH degrees}
The entries specify exactly $Z_n=\sum_wc_w\,L(w)$ for $n=7,8$ in the standard
Lyndon bracketing \eqref{eq:lyndon}; there are no omitted nonzero terms in
this convention. The CSV file continues these expansions through degree
twelve.

\begingroup\small
\renewcommand{\arraystretch}{1.3}
\begin{longtable}{@{}lr@{\hspace{3em}}lr@{}}
\caption{Complete BCH homogeneous polynomials of degrees seven and eight in Lyndon bracketing.}\label{tab:bch78}\\
\toprule
\multicolumn{2}{c}{$Z_7$} & \multicolumn{2}{c}{$Z_8$} \\
$w$ & coefficient of $L(w)$ & $w$ & coefficient of $L(w)$ \\
\midrule\endfirsthead
\multicolumn{4}{c}{\tablename\ \thetable\ (continued)}\\
\toprule
\multicolumn{2}{c}{$Z_7$} & \multicolumn{2}{c}{$Z_8$} \\
$w$ & coefficient of $L(w)$ & $w$ & coefficient of $L(w)$ \\
\midrule\endhead
\bottomrule\endfoot
\texttt{XXXXXXY} & $\tfrac{1}{30240}$ & \texttt{XXXXXXYY} & $\tfrac{1}{60480}$ \\
\texttt{XXXXXYY} & $-\tfrac{1}{5040}$ & \texttt{XXXXXYXY} & $-\tfrac{1}{15120}$ \\
\texttt{XXXXYXY} & $\tfrac{1}{10080}$ & \texttt{XXXXXYYY} & $-\tfrac{1}{10080}$ \\
\texttt{XXXXYYY} & $\tfrac{1}{3780}$ & \texttt{XXXXYXXY} & $\tfrac{1}{20160}$ \\
\texttt{XXXYXXY} & $\tfrac{1}{10080}$ & \texttt{XXXXYXYY} & $-\tfrac{1}{20160}$ \\
\texttt{XXXYXYY} & $\tfrac{1}{1680}$ & \texttt{XXXXYYXY} & $\tfrac{1}{2520}$ \\
\texttt{XXXYYXY} & $\tfrac{1}{1260}$ & \texttt{XXXXYYYY} & $\tfrac{23}{120960}$ \\
\texttt{XXXYYYY} & $\tfrac{1}{3780}$ & \texttt{XXXYXXYY} & $\tfrac{1}{4032}$ \\
\texttt{XXYXXYY} & $\tfrac{1}{2016}$ & \texttt{XXXYXYXY} & $-\tfrac{1}{10080}$ \\
\texttt{XXYXYXY} & $-\tfrac{1}{5040}$ & \texttt{XXXYXYYY} & $\tfrac{13}{30240}$ \\
\texttt{XXYXYYY} & $\tfrac{13}{15120}$ & \texttt{XXXYYXYY} & $\tfrac{1}{20160}$ \\
\texttt{XXYYXYY} & $\tfrac{1}{10080}$ & \texttt{XXXYYYXY} & $-\tfrac{1}{3024}$ \\
\texttt{XXYYYXY} & $-\tfrac{1}{1512}$ & \texttt{XXXYYYYY} & $-\tfrac{1}{10080}$ \\
\texttt{XXYYYYY} & $-\tfrac{1}{5040}$ & \texttt{XXYXYXYY} & $\tfrac{1}{2520}$ \\
\texttt{XYXYXYY} & $\tfrac{1}{1260}$ & \texttt{XXYXYYYY} & $-\tfrac{1}{4032}$ \\
\texttt{XYXYYYY} & $-\tfrac{1}{2016}$ & \texttt{XXYYXYYY} & $-\tfrac{1}{10080}$ \\
\texttt{XYYXYYY} & $-\tfrac{1}{5040}$ & \texttt{XXYYYYYY} & $\tfrac{1}{60480}$ \\
\texttt{XYYYYYY} & $\tfrac{1}{30240}$ & -- & -- \\
\end{longtable}
\endgroup

\subsection{Complete fifth and sixth Zassenhaus exponents}
These are the next two complete exponents after the terms displayed on the
source page. A coefficient multiplies $L(w)$, not $R(w)$. Every entry follows
from \eqref{eq:residual}, from \eqref{eq:frec}, and from
\eqref{eq:treeformula}; the three constructions agree exactly.

\begingroup\small
\renewcommand{\arraystretch}{1.3}
\begin{longtable}{@{}lr@{\hspace{3em}}lr@{}}
\caption{Complete fifth and sixth Zassenhaus exponents in Lyndon bracketing.}\label{tab:zass56}\\
\toprule
\multicolumn{2}{c}{$C_5$} & \multicolumn{2}{c}{$C_6$} \\
$w$ & coefficient of $L(w)$ & $w$ & coefficient of $L(w)$ \\
\midrule\endfirsthead
\multicolumn{4}{c}{\tablename\ \thetable\ (continued)}\\
\toprule
\multicolumn{2}{c}{$C_5$} & \multicolumn{2}{c}{$C_6$} \\
$w$ & coefficient of $L(w)$ & $w$ & coefficient of $L(w)$ \\
\midrule\endhead
\bottomrule\endfoot
\texttt{XXXXY} & $\tfrac{1}{120}$ & \texttt{XXXXXY} & $-\tfrac{1}{720}$ \\
\texttt{XXXYY} & $-\tfrac{1}{30}$ & \texttt{XXXXYY} & $\tfrac{1}{144}$ \\
\texttt{XXYXY} & $-\tfrac{1}{60}$ & \texttt{XXXYYY} & $-\tfrac{1}{72}$ \\
\texttt{XXYYY} & $\tfrac{1}{20}$ & \texttt{XXYYYY} & $\tfrac{1}{72}$ \\
\texttt{XYXYY} & $-\tfrac{1}{20}$ & \texttt{XYXYYY} & $-\tfrac{1}{72}$ \\
\texttt{XYYYY} & $-\tfrac{1}{30}$ & \texttt{XYYYYY} & $-\tfrac{1}{144}$ \\
\end{longtable}
\endgroup

\subsection{Tree polynomials for the Zassenhaus exponents}
Each $C_n$ below is written as a polynomial in the ordered raw polynomials
$A_j$ of \eqref{eq:An}, obtained by summing the weights $\omega(T)$ of
\cref{thm:tree} over trees with the same leaf sequence. Products are
noncommutative and are read from left to right; $A_2^2A_3$ and $A_3A_2^2$ are
different polynomials.
\begin{center}\small
\renewcommand{\arraystretch}{1.3}
\begin{tabular}{@{}r>{\raggedright\arraybackslash}p{0.9\textwidth}@{}}
\toprule
$n$ & $C_n$ in terms of the $A_j$\\
\midrule
2 & $A_{2}$\\
3 & $A_{3}$\\
4 & $A_{4}$ $-\tfrac{1}{2}A_{2}^{2}$\\
5 & $A_{5}$ $-A_{2}A_{3}$\\
6 & $A_{6}$ $+\tfrac{1}{3}A_{2}^{3}$ $-A_{2}A_{4}$ $-\tfrac{1}{2}A_{3}^{2}$\\
7 & $A_{7}$ $+\tfrac{1}{2}A_{2}^{2}A_{3}$ $-A_{2}A_{5}$ $-A_{3}A_{4}$ $+\tfrac{1}{2}A_{3}A_{2}^{2}$\\
8 & $A_{8}$ $-\tfrac{1}{4}A_{2}^{4}$ $+\tfrac{3}{4}A_{2}^{2}A_{4}$ $-A_{2}A_{6}$ $-A_{3}A_{5}$ $+A_{3}A_{2}A_{3}$ $-\tfrac{1}{2}A_{4}^{2}$ $+\tfrac{1}{4}A_{4}A_{2}^{2}$\\
9 & $A_{9}$ $-\tfrac{2}{3}A_{2}^{3}A_{3}$ $+A_{2}^{2}A_{5}$ $-A_{2}A_{7}$ $+\tfrac{1}{3}A_{3}^{3}$ $-A_{3}A_{6}$ $-\tfrac{1}{3}A_{3}A_{2}^{3}$ $+A_{3}A_{2}A_{4}$ $-A_{4}A_{5}$ $+A_{4}A_{2}A_{3}$\\
10 & $A_{10}$ $+\tfrac{1}{5}A_{2}^{5}$ $-\tfrac{2}{3}A_{2}^{3}A_{4}$ $-\tfrac{1}{4}A_{2}^{2}A_{3}^{2}$ $+A_{2}^{2}A_{6}$ $+\tfrac{1}{2}A_{2}A_{3}A_{5}$ $-\tfrac{1}{2}A_{2}A_{3}A_{2}A_{3}$ $-A_{2}A_{8}$ $+\tfrac{1}{2}A_{3}^{2}A_{4}$ $-\tfrac{1}{4}A_{3}^{2}A_{2}^{2}$ $-A_{3}A_{7}$ $-\tfrac{1}{2}A_{3}A_{2}^{2}A_{3}$ $+A_{3}A_{2}A_{5}$ $-A_{4}A_{6}$ $-\tfrac{1}{3}A_{4}A_{2}^{3}$ $+A_{4}A_{2}A_{4}$ $+\tfrac{1}{2}A_{4}A_{3}^{2}$ $-\tfrac{1}{2}A_{5}^{2}$ $+\tfrac{1}{2}A_{5}A_{2}A_{3}$\\
\bottomrule
\end{tabular}
\end{center}

\subsection{Sizes of the exported expansions}
\begin{table}[htbp]
\centering\small
\begin{tabular}{@{}rrrrr@{}}
\toprule
Degree & BCH words & BCH Lyndon entries & $C_n$ words & $C_n$ Lyndon entries\\
\midrule
1 & 2 & 2 & 0 & 0 \\
2 & 2 & 1 & 2 & 1 \\
3 & 6 & 2 & 6 & 2 \\
4 & 4 & 1 & 12 & 3 \\
5 & 30 & 6 & 30 & 6 \\
6 & 28 & 5 & 42 & 6 \\
7 & 126 & 18 & 126 & 18 \\
8 & 124 & 17 & 234 & 29 \\
9 & 390 & 55 & 510 & 56 \\
10 & 388 & 55 & 954 & 95 \\
11 & 2046 & 186 & 2046 & 186 \\
12 & 2044 & 185 & 3986 & 326 \\
\bottomrule
\end{tabular}
\caption{Numbers of nonzero coefficients in the supplied encodings. The
degree-one Zassenhaus factor $e^{Y}$ is separated from the exponents $C_n$,
$n\geq2$.}\label{tab:counts}
\end{table}
